\documentclass[11pt,a4paper]{article}

\usepackage[utf8]{inputenc}
\usepackage[T1]{fontenc}
\usepackage{lmodern}
\usepackage[spanish,es-tabla,es-noshorthands]{babel}
\usepackage[a4paper,margin=2.2cm]{geometry}
\usepackage{xcolor}
\usepackage{listings}
\usepackage{tcolorbox}
\usepackage{enumitem}
\usepackage{titlesec}
\usepackage{parskip}
\usepackage{hyperref}
\usepackage{tikz}
\usetikzlibrary{arrows.meta,positioning,shapes.geometric}

\definecolor{azulOscuro}{HTML}{1E3C78}
\definecolor{azulMedio}{HTML}{2A5BB1}
\definecolor{azulSuave}{HTML}{E8EEFA}
\definecolor{grisCodigo}{HTML}{F5F6FA}
\definecolor{grisBorde}{HTML}{D0D4DD}
\definecolor{kw}{HTML}{0033B3}
\definecolor{str}{HTML}{067D17}
\definecolor{com}{HTML}{8C8C8C}
\definecolor{verdeTAC}{HTML}{2D7D46}
\definecolor{rojoOpt}{HTML}{B23B3B}
\definecolor{naranjaPausa}{HTML}{E08600}
\definecolor{moradoSem}{HTML}{8E44AD}

\hypersetup{
  colorlinks=true,
  linkcolor=azulMedio,
  urlcolor=azulMedio,
  pdftitle={Compiladores 2026 - Guia Fase 2}
}

\titleformat{\section}
  {\color{azulOscuro}\Large\bfseries}
  {\thesection.}{0.6em}{}
\titleformat{\subsection}
  {\color{azulMedio}\large\bfseries}
  {\thesubsection.}{0.5em}{}
\titleformat{\subsubsection}
  {\color{azulMedio}\normalsize\bfseries}
  {\thesubsubsection.}{0.5em}{}

\lstdefinestyle{py}{
  language=Python,
  backgroundcolor=\color{grisCodigo},
  basicstyle=\ttfamily\footnotesize,
  keywordstyle=\color{kw}\bfseries,
  stringstyle=\color{str},
  commentstyle=\color{com}\itshape,
  frame=single, framerule=0.4pt, rulecolor=\color{grisBorde},
  numbers=none, showstringspaces=false, breaklines=true,
  xleftmargin=8pt, xrightmargin=8pt, aboveskip=6pt, belowskip=6pt,
  inputencoding=utf8, extendedchars=true,
  literate=
    {á}{{\'a}}1 {é}{{\'e}}1 {í}{{\'i}}1 {ó}{{\'o}}1 {ú}{{\'u}}1
    {Á}{{\'A}}1 {É}{{\'E}}1 {Í}{{\'I}}1 {Ó}{{\'O}}1 {Ú}{{\'U}}1
    {ñ}{{\~n}}1 {Ñ}{{\~N}}1 {¡}{{!`}}1 {¿}{{?`}}1,
}
\lstdefinestyle{tac}{
  backgroundcolor=\color{grisCodigo},
  basicstyle=\ttfamily\footnotesize,
  frame=single, framerule=0.4pt, rulecolor=\color{grisBorde},
  numbers=none, showstringspaces=false, breaklines=true,
  xleftmargin=8pt, xrightmargin=8pt, aboveskip=6pt, belowskip=6pt,
}
\lstdefinestyle{ml}{
  backgroundcolor=\color{azulSuave},
  basicstyle=\ttfamily\footnotesize,
  frame=single, framerule=0.4pt, rulecolor=\color{azulMedio},
  numbers=none, showstringspaces=false, breaklines=true,
  xleftmargin=8pt, xrightmargin=8pt, aboveskip=6pt, belowskip=6pt,
}

\newtcolorbox{caja}[1]{
  colback=grisCodigo, colframe=azulMedio,
  title=\textbf{#1}, fonttitle=\color{white},
  boxrule=0.5pt, arc=2pt,
  left=6pt,right=6pt,top=4pt,bottom=4pt,
  before skip=8pt, after skip=8pt,
}
\newtcolorbox{pausa}[1]{
  colback=naranjaPausa!10, colframe=naranjaPausa,
  title=\textbf{#1}, fonttitle=\color{white},
  boxrule=0.7pt, arc=2pt,
  left=6pt,right=6pt,top=4pt,bottom=4pt,
  before skip=8pt, after skip=8pt,
}
\newtcolorbox{concepto}[1]{
  colback=verdeTAC!10, colframe=verdeTAC,
  title=\textbf{#1}, fonttitle=\color{white},
  boxrule=0.6pt, arc=2pt,
  left=6pt,right=6pt,top=4pt,bottom=4pt,
  before skip=8pt, after skip=8pt,
}
\newtcolorbox{rubrica}[1]{
  colback=moradoSem!10, colframe=moradoSem,
  title=\textbf{#1}, fonttitle=\color{white},
  boxrule=0.6pt, arc=2pt,
  left=6pt,right=6pt,top=4pt,bottom=4pt,
  before skip=8pt, after skip=8pt,
}

\title{\vspace{-1.2cm}\color{azulOscuro}\textbf{Compiladores 2026 --- Fase 2}\\[0.3em]
\large\color{black}Frontend/Backend, Codigo de Tres Direcciones, Optimizacion y Reportes HTML}
\author{}
\date{}

\begin{document}
\maketitle
\vspace{-1.2cm}

\section{Mapa de la presentacion}

\begin{enumerate}[leftmargin=*,nosep]
  \item Resumen ejecutivo: que se entrega en esta fase 2.
  \item Arquitectura: separacion \textbf{Frontend / Backend}.
  \item \textbf{PAUSA teorica:} tipos de codigo intermedio y cual usamos.
  \item Codigo de Tres Direcciones (TAC).
  \item Pipeline PLY: \texttt{lex} + \texttt{yacc} produciendo el AST.
  \item Generacion del TAC a partir del AST.
  \item Analisis semantico --- \texttt{check\_semantic}.
  \item \textbf{PAUSA pedagogica:} visualizacion y reportes HTML.
  \item Optimizacion del TAC.
  \item Pruebas y QA contra la rubrica.
\end{enumerate}

\section{Resumen ejecutivo}

\begin{caja}{Que entregamos en la fase 2}
\begin{itemize}[leftmargin=*,nosep]
  \item Reestructuracion del codigo en \texttt{frontend/} y \texttt{backend/}
        con un contrato compartido (\texttt{intermedio.py}).
  \item \textbf{Generador de codigo de tres direcciones (TAC)} sobre el
        AST que produce el parser PLY yacc.
  \item \textbf{Analisis semantico} (\texttt{check\_semantic}) que detecta
        tipos incompatibles, variables no declaradas, condiciones invalidas,
        division por cero, etc., con \emph{linea y columna} por error.
  \item \textbf{Optimizador} con cinco pasadas iterando hasta punto fijo.
  \item \textbf{Reportes HTML} estructurados de tokens, tabla de simbolos,
        errores (con separado para semanticos), TAC original y TAC optimizado.
  \item UI dual: escritorio (\texttt{Tkinter}, \texttt{main.py}) y web
        (\texttt{index.php} + \texttt{analizar.php}) que muestran
        todo el pipeline al usuario.
  \item Suite de QA (\texttt{qa\_rubrica.py}) con 42 casos automaticos
        cubriendo los cinco criterios de la rubrica.
\end{itemize}
\end{caja}

\section{Arquitectura: Frontend / Backend}

\subsection{El principio}

Un compilador moderno separa lo que depende del \emph{lenguaje fuente}
(frontend) de lo que depende de la \emph{maquina destino} (backend). Entre
los dos esta la \textbf{representacion intermedia} (IR), que es el
contrato comun.

\begin{center}
\begin{tikzpicture}[node distance=1.4cm and 1.0cm,
  every node/.style={font=\small\sffamily,align=center},
  caja/.style={draw=azulMedio,fill=azulSuave,thick,rounded corners=2pt,
               minimum width=2.6cm,minimum height=1.0cm},
  arr/.style={-{Stealth[length=2.5mm]},thick,azulMedio}]
\node[caja] (fe) {FRONTEND\\\scriptsize lenguaje fuente};
\node[caja, right=of fe] (ir) {IR\\\scriptsize TAC};
\node[caja, right=of ir] (be) {BACKEND\\\scriptsize maquina destino};
\draw[arr] (fe) -- (ir);
\draw[arr] (ir) -- (be);
\node[below=0.1cm of fe,font=\scriptsize] {PLY lex + yacc};
\node[below=0.1cm of ir,font=\scriptsize] {Quad (op,a1,a2,d)};
\node[below=0.1cm of be,font=\scriptsize] {Optimizador};
\end{tikzpicture}
\end{center}

\subsection{Como lo aplicamos}

\begin{lstlisting}[style=tac]
compiladores-2026/
|-- frontend/                       <- depende del lenguaje fuente
|   |-- __init__.py                 reexporta API publica
|   |-- lexer.py                    PLY: tokens + tabla de simbolos
|   |-- tabla_simbolos.py           Simbolo + TablaSimbolos con ambitos
|   |-- errores.py                  agregar_lexico/sintactico/semantico
|   |-- parser.py                   PLY yacc: AST + check_semantic
|   `-- generador_intermedio.py     AST -> TAC
|
|-- intermedio.py                   <- contrato compartido (IR)
|   `-- Quad(op, arg1, arg2, dest), formatear_tac(...)
|
|-- backend/                        <- depende de la maquina objetivo
|   |-- __init__.py
|   `-- optimizador.py              5 pasadas hasta punto fijo
|
|-- reportes.py                     reportes HTML (semanticos, TAC, etc.)
|-- bridge.py                       pipeline -> JSON
|-- main.py                         UI Tkinter
|-- index.php / analizar.php        UI web
\end{lstlisting}

\begin{lstlisting}[style=py]
# frontend/__init__.py
from .lexer            import Lexer
from .tabla_simbolos   import TablaSimbolos, Simbolo
from .parser           import build_tree, get_kids, node_label, \
                              TIPOS_DATO, RESERVED, check_semantic
from .generador_intermedio import GeneradorTAC

# backend/__init__.py
from .optimizador import optimizar
\end{lstlisting}

\textbf{Punto de discusion:} \texttt{intermedio.py} no importa de
\texttt{frontend} ni de \texttt{backend}. Es la \emph{frontera}: cualquier
modulo que respete \texttt{Quad} puede leer o escribir TAC.

\section{PAUSA: tipos de codigo intermedio}

\begin{pausa}{Pausa teorica en la presentacion}
Antes de mostrar nuestro generador de TAC conviene explicar \textbf{que
tipos de codigo intermedio existen} y \textbf{cual elegimos}.
\end{pausa}

\subsection{Los tres tipos clasicos}

\begin{concepto}{1. Notacion postfija (Reverse Polish)}
La operacion se escribe \emph{despues} de los operandos. \\
\textbf{Ejemplo:} \texttt{a + b * c} se vuelve \texttt{a b c * +}. \\
\textbf{Ventaja:} compacta y facil de evaluar con una pila. \\
\textbf{Desventaja:} no representa flujo de control (saltos, etiquetas),
solo expresiones aritmeticas.
\end{concepto}

\begin{concepto}{2. Arbol Sintactico Abstracto (AST)}
Estructura de arbol: cada nodo es una operacion o un literal y los hijos
son los operandos. \\
\textbf{Ejemplo:} \texttt{a + b * c} es un nodo \texttt{+} cuyos hijos
son \texttt{a} y un sub-nodo \texttt{*} con hijos \texttt{b} y \texttt{c}. \\
\textbf{Ventaja:} mantiene la estructura del programa, ideal para
analisis semantico. \\
\textbf{Desventaja:} no es lineal, no se ejecuta secuencialmente sin
recorrerlo.
\end{concepto}

\begin{concepto}{3. Codigo de Tres Direcciones (TAC)}
Cada instruccion tiene a lo mas tres operandos: dos fuentes y un destino.
Forma canonica: \texttt{dest = arg1 op arg2}. \\
\textbf{Ejemplo:} \texttt{a + b * c} se descompone en dos cuadruplos. \\
\textbf{Ventaja:} lineal (una instruccion por linea), explicito
(temporales con nombre), parecido al ensamblador. \\
\textbf{Desventaja:} mas voluminoso que la postfija.
\end{concepto}

\begin{lstlisting}[style=tac]
$t1 = b * c
$t2 = a + $t1
\end{lstlisting}

\subsection{Cual usa nuestro compilador?}

\begin{caja}{Decision de diseno}
El compilador produce y usa internamente \textbf{dos representaciones}:
\begin{enumerate}[leftmargin=*,nosep]
  \item \textbf{AST} producido por \texttt{build\_tree()} en
        \texttt{frontend/parser.py}. Lo usamos para analisis semantico,
        para visualizar el arbol y como entrada del generador de TAC.
  \item \textbf{Codigo de Tres Direcciones (TAC)} producido por
        \texttt{GeneradorTAC.generar()} en
        \texttt{frontend/generador\_intermedio.py}. Es el formato sobre
        el que trabaja el optimizador (\texttt{backend/optimizador.py}).
\end{enumerate}
\textbf{No usamos notacion postfija} porque no permite expresar saltos
ni asignaciones, que son basicos en este lenguaje.
\end{caja}

\section{Codigo de Tres Direcciones}

\subsection{Estructura del cuadruplo}

\begin{lstlisting}[style=py]
# intermedio.py
@dataclass
class Quad:
    op:   str                       # operacion o palabra clave
    arg1: Optional[str] = None      # primer operando
    arg2: Optional[str] = None      # segundo operando
    dest: Optional[str] = None      # destino o etiqueta destino
\end{lstlisting}

\subsection{Conjunto de instrucciones del IR}

\begin{lstlisting}[style=tac]
Asignacion:        dest = arg1                     op = "="
Operacion binaria: dest = arg1 op arg2             op in {+,-,*,/,%,==,!=,<,>,<=,>=,&&,||}
Operacion unaria:  dest = ! arg1                   op = "!"
Etiqueta:          L:                              op = "label", dest = L
Salto:             goto L                          op = "goto", dest = L
Salto condicional: ifFalse arg1 goto L             op = "if_false", dest = L
Imprimir:          imprimir arg1                   op = "print"
Leer:              leer dest                       op = "read", dest = nombre
\end{lstlisting}

\subsection{Pretty-printer \texttt{formatear\_tac}}

\begin{lstlisting}[style=py]
def formatear_tac(quads):
    filas, n = [], 0
    for q in quads:
        op, a1, a2, d = q.op, q.arg1, q.arg2, q.dest
        if   op == "label":     instr = f"{d}:"
        elif op == "goto":      instr = f"goto {d}"
        elif op == "if_false":  instr = f"ifFalse {a1} goto {d}"
        elif op == "=":         instr = f"{d} = {a1}"
        elif op == "print":     instr = f"imprimir {a1}"
        elif op == "read":      instr = f"leer {d}"
        elif op == "!":         instr = f"{d} = !{a1}"
        else:                   instr = f"{d} = {a1} {op} {a2}"
        n += 1
        filas.append({"n": n, "instruccion": instr,
                      "op": op, "arg1": a1 or "-",
                      "arg2": a2 or "-", "dest": d or "-"})
    return filas
\end{lstlisting}

\section{Pipeline PLY: lex + yacc}

\subsection{Lexer con PLY}

\begin{lstlisting}[style=py]
# frontend/lexer.py
import ply.lex as lex
from . import errores as _errores_mod

reservadas = {
    "programa": "PROGRAMA",
    "entero":   "ENTERO",   "decimal":  "DECIMAL",
    "cadena":   "CADENA_TIPO", "booleano": "BOOLEANO",
    "si": "SI", "sino": "SINO",
    "mientras": "MIENTRAS", "hacer_mientras": "HACER_MIENTRAS",
    "para": "PARA", "imprimir": "IMPRIMIR", "leer": "LEER",
    "verdadero": "VERDADERO", "falso": "FALSO",
    "funcion": "FUNCION", "procedimiento": "PROCEDIMIENTO",
    "retornar": "RETORNAR",
}

t_MAS = r'\+';  t_MENOS = r'-';  t_MULTIPLICACION = r'\*';
t_DIVIDIR = r'/';  t_MODULO = r'%';
t_IGUAL = r'==';  t_DIFERENTE = r'!='
t_MENOR_IGUAL = r'<=';  t_MAYOR_IGUAL = r'>='
t_AND = r'&&';  t_OR = r'\|\|';  t_NOT = r'!'
...

def t_IDENTIFICADOR(t):
    r'[a-zA-Z_][a-zA-Z0-9_]*'
    t.type = reservadas.get(t.value, 'IDENTIFICADOR')
    return t

def t_error(t):
    col = encontrar_columna(t.lexer.lexdata, t)
    _errores_mod.agregar_lexico(
        f"Caracter ilegal '{t.value[0]}'",
        t.lineno, col, valor=t.value[0],
    )
    t.lexer.skip(1)
\end{lstlisting}

\textbf{Que hace:}
\begin{itemize}[leftmargin=*,nosep]
  \item Define tokens con regex como atributos \texttt{t\_XXX} (PLY los
        recoge por convencion de nombre).
  \item Las palabras reservadas se distinguen del identificador
        \emph{post-match}: si el \texttt{IDENTIFICADOR} matchea
        \texttt{"programa"}, se reclasifica a \texttt{PROGRAMA}.
  \item Los errores lexicos se acumulan en el modulo \texttt{errores}
        (con linea y columna).
\end{itemize}

\subsection{Parser con PLY yacc}

\begin{lstlisting}[style=py]
# frontend/parser.py
import ply.yacc as yacc

precedence = (
    ('left',  'OR'),
    ('left',  'AND'),
    ('left',  'IGUAL', 'DIFERENTE'),
    ('left',  'MENOR', 'MAYOR', 'MENOR_IGUAL', 'MAYOR_IGUAL'),
    ('left',  'MAS', 'MENOS'),
    ('left',  'MULTIPLICACION', 'DIVIDIR', 'MODULO'),
    ('right', 'NOT', 'UMENOS'),
)

def p_programa(p):
    '''programa : PROGRAMA LLAVE_IZQ sentencias LLAVE_DER
                | sentencias'''
    if len(p) == 5:
        p[0] = {"type": "Program", "children": p[3]}
    else:
        p[0] = {"type": "Program", "children": p[1]}

def p_declaracion(p):
    '''declaracion : tipo IDENTIFICADOR ASIGNAR expresion PUNTO_COMA'''
    p[0] = {
        "type":     "Declaration",
        "dataType": p[1],
        "id":       _tok_de_p(p, 2, tipo="IDENTIFICADOR"),
        "expr":     p[4],
    }

def p_expresion_binaria(p):
    '''expresion : expresion MAS expresion
                 | expresion MENOS expresion
                 ...
                 | expresion AND expresion
                 | expresion OR expresion'''
    op_tok = {"tipo": p.slice[2].type,
              "valor": _OP_TIPOS.get(p.slice[2].type, str(p[2])),
              "linea": p.lineno(2), "columna": col}
    p[0] = {"type": "BinaryOp", "op": op_tok, "left": p[1], "right": p[3]}
\end{lstlisting}

\textbf{Decision clave:} las acciones semanticas del parser \emph{no
producen tuplas} (como hacia la fase 1) sino \textbf{dicts en formato AST}
que el resto del pipeline (visitor del TAC, semantico, visualizacion del
arbol) ya entiende.

\subsection{Sintaxis del lenguaje}

\begin{lstlisting}[style=ml]
programa {
    entero edad = 25;
    decimal salario = 15750.50;
    cadena nombre = "Ana Garcia";
    booleano activo = verdadero;

    si (edad >= 18 && activo) {
        imprimir(nombre);
    } sino {
        imprimir("menor");
    }

    mientras (edad > 0) { edad = edad - 1; }
    para (entero i = 0; i < 5; i = i + 1) { imprimir(i); }
}
\end{lstlisting}

\section{Generacion del TAC}

\subsection{Visitor por tipo de nodo}

\begin{lstlisting}[style=py]
class GeneradorTAC:
    def __init__(self):
        self.codigo: list[Quad] = []
        self._t = 0   # contador de temporales
        self._l = 0   # contador de etiquetas

    def _nuevo_temp(self):  self._t += 1; return f"$t{self._t}"
    def _nueva_etiq(self):  self._l += 1; return f"$L{self._l}"
    def emit(self, op, a1=None, a2=None, d=None):
        self.codigo.append(Quad(op, a1, a2, d))

    def generar(self, ast):
        self.codigo = []; self._t = 0; self._l = 0
        if ast: self._gen(ast)
        return self.codigo

    def _gen(self, node):
        if not node: return None
        method = getattr(self, f"_g_{node['type']}", self._g_default)
        return method(node)
\end{lstlisting}

\textbf{Patron visitor:} para un nodo \texttt{If} se llama
\texttt{\_g\_If}, para \texttt{BinaryOp} se llama \texttt{\_g\_BinaryOp},
y asi. Sin \texttt{if/elif} repetitivos.

\subsection{Expresiones binarias en post-orden}

\begin{lstlisting}[style=py]
if t == "BinaryOp":
    izq = self._g_expr(node["left"])
    der = self._g_expr(node["right"])
    tmp = self._nuevo_temp()
    op  = _BIN_TIPOS.get(node["op"]["tipo"], node["op"]["valor"])
    self.emit(op, izq, der, tmp)
    return tmp
\end{lstlisting}

\textbf{Como funciona:} primero se genera el TAC del subarbol izquierdo
y se obtiene el lugar donde quedo su resultado; lo mismo con el derecho;
despues se emite \texttt{tmp = izq op der}. Los temporales se encadenan
naturalmente.

\subsection{Estructuras de control}

\textbf{Si / sino:}
\begin{lstlisting}[style=py]
def _g_If(self, node):
    cond   = self._g_expr(node["cond"])
    L_else = self._nueva_etiq()
    L_end  = self._nueva_etiq() if node.get("elseBody") else L_else
    self.emit("if_false", cond, None, L_else)
    self._gen(node["body"])
    if node.get("elseBody"):
        self.emit("goto",  None, None, L_end)
        self.emit("label", None, None, L_else)
        self._gen(node["elseBody"])
        self.emit("label", None, None, L_end)
    else:
        self.emit("label", None, None, L_else)
\end{lstlisting}

\textbf{Mientras:}
\begin{lstlisting}[style=py]
def _g_While(self, node):
    L_inicio = self._nueva_etiq()
    L_fin    = self._nueva_etiq()
    self.emit("label",    None, None, L_inicio)
    cond = self._g_expr(node["cond"])
    self.emit("if_false", cond, None, L_fin)
    self._gen(node["body"])
    self.emit("goto",     None, None, L_inicio)
    self.emit("label",    None, None, L_fin)
\end{lstlisting}

\subsection{Ejemplo end-to-end}

\textbf{Fuente:}
\begin{lstlisting}[style=ml]
programa {
    entero a = 3;
    entero b = 4;
    entero c = (a + b) * 2;
    si (c > 10) { imprimir(c); } sino { imprimir("pequeno"); }
}
\end{lstlisting}

\textbf{TAC generado (12 cuadruplos):}
\begin{lstlisting}[style=tac]
 1: a = 3
 2: b = 4
 3: $t1 = a + b
 4: $t2 = $t1 * 2
 5: c = $t2
 6: $t3 = c > 10
 7: ifFalse $t3 goto $L1
 8: imprimir c
 9: goto $L2
10: $L1:
11: imprimir "pequeno"
12: $L2:
\end{lstlisting}

\textbf{TAC optimizado (7 cuadruplos --- reduccion 41.7\%):}
\begin{lstlisting}[style=tac]
 1: a = 3
 2: b = 4
 3: c = 14
 4: imprimir 14
 5: goto $L2
 6: imprimir "pequeno"
 7: $L2:
\end{lstlisting}

\textbf{Punto de discusion:} la propagacion de constantes y el folding
calculan \texttt{(3+4)*2 = 14} en tiempo de compilacion. Los temporales
muertos se eliminan. La \emph{rama muerta} (linea 6) se conserva porque
no podemos garantizar staticamente que ese imprimir no sera alcanzado
desde otro flujo, pero los temporales asociados ya estan plegados.

\section{Analisis semantico --- \texttt{check\_semantic}}

\begin{rubrica}{Criterio 1 de la rubrica (25 pts)}
La rubrica exige detectar errores semanticos. Nuestro
\texttt{check\_semantic} recorre el AST e identifica:
\begin{itemize}[leftmargin=*,nosep]
  \item Variables \textbf{no declaradas} en uso, asignacion o lectura.
  \item Variables \textbf{duplicadas} en el mismo ambito (lo detecta la
        tabla de simbolos al insertar).
  \item \textbf{Tipos incompatibles} en declaracion o asignacion
        (\texttt{entero = "hola"}, \texttt{booleano = 5}\dots).
  \item \textbf{Condiciones} de \texttt{si}/\texttt{mientras}/
        \texttt{hacer\_mientras}/\texttt{para} que no son booleanas.
  \item \textbf{Division por cero} con denominador literal.
  \item \textbf{Operacion aritmetica} entre tipos incompatibles
        (e.g.\ \texttt{booleano + entero}).
  \item \textbf{Comparaciones} y \textbf{operaciones logicas} entre
        tipos invalidos.
  \item \textbf{Operadores unarios} \texttt{-} y \texttt{!} aplicados a
        tipo erroneo.
  \item Funcion no definida en llamadas.
  \item \emph{Widening} permitido: \texttt{decimal = entero} se acepta
        sin error (promocion entero $\rightarrow$ decimal).
\end{itemize}
\textbf{Cada error reporta linea y columna.}
\end{rubrica}

\subsection{Compatibilidad de tipos}

\begin{lstlisting}[style=py]
def _tipo_compatible(declarado, expresion):
    if declarado == expresion:
        return True
    if declarado == "decimal" and expresion == "entero":
        return True   # widening entero -> decimal
    return False
\end{lstlisting}

\subsection{Inferencia de tipos en expresiones}

\begin{lstlisting}[style=py]
def get_type(node):
    if t == "Literal":   return _tipo_literal(node["token"])
    if t == "StringLit": return "cadena"
    if t == "BoolLit":   return "booleano"
    if t == "Identifier":
        sym = tabla.buscar(node["token"]["valor"])
        if not sym:
            _sem("Variable '...' no declarada", linea, col)
            return None
        return sym.tipo
    if t == "BinaryOp":
        lt = get_type(node["left"]); rt = get_type(node["right"])
        op = node["op"]["tipo"]
        # Division por cero
        if op in ("DIVIDIR", "MODULO") and node["right"].get("type") == "Literal":
            if float(node["right"]["token"]["valor"]) == 0:
                _sem("Division por cero", op_lin, op_col)
        # Aritmetica entre numericos
        if op in ("MAS","MENOS","MULTIPLICACION","DIVIDIR","MODULO"):
            if lt == "entero" and rt == "entero":  return "entero"
            if lt in ("entero","decimal") and rt in ("entero","decimal"):
                return "decimal"
            _sem(f"Operacion aritmetica incompatible entre '{lt}' y '{rt}'",
                 op_lin, op_col)
            return None
        ...
\end{lstlisting}

\subsection{Reglas por sentencia}

\begin{lstlisting}[style=py]
def check_node(node):
    if t == "Declaration":
        et = get_type(node["expr"])
        vt = _TIPO_VAR_DE_TOKEN[node["dataType"]["tipo"]]
        if et and not _tipo_compatible(vt, et):
            _sem(f"Asignacion incompatible: '{vt}' no puede recibir '{et}'",
                 linea, col)
    elif t == "Assignment":
        sym = tabla.buscar(node["id"]["valor"])
        if not sym:
            _sem(f"Variable no declarada", linea, col)
        else:
            et = get_type(node["expr"])
            if et and not _tipo_compatible(sym.tipo, et):
                _sem(f"Asignacion incompatible", linea, col)
    elif t == "If":
        ct = get_type(node["cond"])
        if ct and ct != "booleano":
            _sem(f"Condicion del 'si' debe ser booleana, no '{ct}'", ...)
        check_node(node["body"])
        ...
\end{lstlisting}

\section{PAUSA: visualizacion y reportes HTML}

\begin{pausa}{Pausa pedagogica}
Aqui se hace pausa para mostrar como el compilador \emph{exhibe} su
trabajo. Esto cumple dos criterios de la rubrica: el reporte HTML
de errores semanticos (criterio 2, 20 pts) y el funcionamiento general
del compilador (criterio 5, 10 pts).
\end{pausa}

\subsection{UI Tkinter (\texttt{main.py})}

Una sola clase \texttt{MiniIDE(tk.Tk)} con seis pestanias:
\begin{itemize}[leftmargin=*,nosep]
  \item \textbf{Editor} con coloreado de tokens en vivo.
  \item \textbf{Tokens} con tipo, valor, linea y columna.
  \item \textbf{Arbol} (AST) renderizado en canvas.
  \item \textbf{Semantico} con tabla de simbolos.
  \item \textbf{Codigo Intermedio} con TAC original, TAC optimizado y metricas.
  \item \textbf{Errores} con consola dedicada.
\end{itemize}

Atajo: \texttt{F5} = Analizar. Boton \textbf{REPORTES HTML} en la barra
inferior que invoca \texttt{reportes.generar\_reportes\_completos(\dots)}.

\subsection{UI web (\texttt{index.php} + \texttt{analizar.php})}

\begin{lstlisting}[style=py]
# analizar.php llama a bridge.py via exec()
$tmpFile = tempnam(sys_get_temp_dir(), 'minicomp_') . '.programa';
file_put_contents($tmpFile, $codigo);
$python  = (PHP_OS_FAMILY === 'Windows') ? 'python' : 'python3';
$cmd     = escapeshellcmd($python).' '
         . escapeshellarg(__DIR__.'/bridge.py').' '
         . escapeshellarg($tmpFile).' 2>&1';
exec($cmd, $output, $returnCode);
@unlink($tmpFile);
echo json_encode(json_decode(implode("\n",$output),true), JSON_UNESCAPED_UNICODE);
\end{lstlisting}

\subsection{\texttt{reportes.py}: HTML estructurado}

\begin{rubrica}{Criterio 2 de la rubrica (20 pts) --- el HTML semantico}
El archivo \texttt{reportes.py} genera HTMLs profesionales con CSS
moderno y consistente. La rubrica pide especificamente un reporte
\emph{estructurado, claro y funcional con linea y columna de cada error
semantico}: ese es \texttt{generar\_html\_errores\_semanticos()}.
\end{rubrica}

\textbf{Funciones expuestas:}

\begin{lstlisting}[style=py]
# reportes.py
generar_html_errores_semanticos(errores, ruta)   # << el clave (rubrica)
generar_html_errores(errores, ruta)              # todos los errores
generar_html_tokens(tokens, ruta)
generar_html_tabla_simbolos(tabla, ruta)
generar_html_tac(quads_formateados, ruta, titulo)
generar_reportes_completos(directorio, ...)      # todo de una
\end{lstlisting}

\textbf{Lo que muestra el reporte semantico:}
\begin{itemize}[leftmargin=*,nosep]
  \item Encabezado con fecha y resumen estadistico por categoria
        (variables no declaradas, tipos incompatibles, division por
        cero, condiciones no booleanas, etc.).
  \item Tabla principal con columnas: \texttt{\#}, tipo (tag
        coloreado), \textbf{linea}, \textbf{columna}, identificador
        afectado, descripcion.
  \item Mensaje verde de "sin errores" si la lista esta vacia.
\end{itemize}

\textbf{Estructura del HTML} (extracto del CSS):

\begin{lstlisting}[style=tac]
header { gradient azul + titulo + fecha }
.summary { chips estadisticos por categoria }
table {
    thead   th  bg=primario, fg=blanco
    tbody   alternancia de filas, hover, monospace para identificadores
    .tag    coloreado por tipo (lex=rojo, sint=amarillo, sem=morado)
}
\end{lstlisting}

\textbf{Como se genera:}

\begin{itemize}[leftmargin=*,nosep]
  \item Desde Tkinter: boton "REPORTES HTML" $\rightarrow$ abre el HTML
        de errores semanticos en el navegador.
  \item Desde CLI: \texttt{python bridge.py archivo.programa --reportes salida/}
        crea todos los archivos en \texttt{salida/}.
\end{itemize}

\section{Optimizacion del TAC}

\subsection{Driver con punto fijo}

\begin{lstlisting}[style=py]
def optimizar(codigo, max_iter: int = 12):
    actual = deepcopy(codigo); traza = []
    pasadas = [
        ("Constant Folding & Algebraic", _constant_folding),
        ("Constant / Copy Propagation",  _propagation),
        ("Branch Pruning",               _branch_pruning),
        ("Dead-Code Elimination",        _dead_code),
        ("Jump Threading",               _jump_threading),
    ]
    for it in range(max_iter):
        cambio_global = False
        for nombre, fn in pasadas:
            antes = len(actual)
            actual, cambio = fn(actual)
            if cambio:
                traza.append({"iter": it+1, "pasada": nombre,
                              "antes": antes, "despues": len(actual),
                              "delta": len(actual) - antes})
                cambio_global = True
        if not cambio_global: break
    return actual, traza
\end{lstlisting}

\textbf{Que hace:} aplica las cinco pasadas \emph{en orden} y repite el
ciclo hasta que una iteracion completa no haga ningun cambio (punto
fijo). Cada modificacion se registra en \texttt{traza}.

\subsection{Pasada 1 --- Constant Folding}

\textbf{Idea:} si los dos operandos son constantes, evaluar en tiempo
de compilacion. Si uno es identidad algebraica (\texttt{x+0},
\texttt{x*1}, \texttt{x*0}), reescribir.

\begin{lstlisting}[style=py]
if q.op == "+" and a1 == "0":
    nuevo.append(Quad("=", a2, None, q.dest));  cambio=True;  continue
if q.op == "*" and (a1 == "0" or a2 == "0"):
    nuevo.append(Quad("=", "0", None, q.dest)); cambio=True;  continue
if _es_numero(a1) and _es_numero(a2):
    r = _BIN_OPS[q.op](_to_num(a1), _to_num(a2))
    nuevo.append(Quad("=", _fmt(r), None, q.dest)); cambio=True
\end{lstlisting}

\subsection{Pasada 2 --- Constant / Copy Propagation}

\textbf{Idea:} mantener un entorno con los valores conocidos. Cuando se
ve \texttt{x = 5}, recordar \texttt{env["x"] = "5"} y reemplazar usos
posteriores. El entorno se invalida al cruzar etiquetas o saltos.

\subsection{Pasada 3 --- Branch Pruning}

\textbf{Idea:} si la propagacion dejo la condicion de un salto como
literal (\texttt{verdadero}/\texttt{falso}), eliminar el salto inutil
o convertirlo en \texttt{goto} incondicional.

\begin{lstlisting}[style=py]
_VERDADERO = frozenset(("true", "True", "1", "verdadero", "Verdadero"))
_FALSO     = frozenset(("false", "False", "0", "falso", "Falso"))

if q.op == "if_false":
    if q.arg1 in _VERDADERO:  cambio=True;  continue   # nunca salta
    if q.arg1 in _FALSO:                                 # siempre salta
        nuevo.append(Quad("goto", None, None, q.dest)); cambio=True
\end{lstlisting}

\subsection{Pasada 4 --- Dead-Code Elimination}

\textbf{Idea:} eliminar los cuadruplos cuyo destino es un temporal que
no se lee en ningun lado.

\begin{lstlisting}[style=py]
leidos = {a for q in codigo for a in (q.arg1, q.arg2) if a}
for q in codigo:
    if q.op == "=" or q.op in _BIN_OPS or q.op == "!":
        if q.dest and _es_temp(q.dest) and q.dest not in leidos:
            cambio=True;  continue   # se descarta
    nuevo.append(q)
\end{lstlisting}

\textbf{Importante:} solo se eliminan temporales (\texttt{\$tN}). Las
variables del usuario nunca se borran aunque parezcan muertas.

\subsection{Pasada 5 --- Jump Threading}

\textbf{Idea (a):} si encontramos \texttt{goto L} seguido de
\texttt{L:}, eliminamos el \texttt{goto}. \\
\textbf{Idea (b):} las etiquetas que ya no son referenciadas se borran.

\subsection{Como cooperan las pasadas}

\begin{caja}{El truco del punto fijo}
\begin{enumerate}[leftmargin=*,nosep]
  \item Folding pliega \texttt{\$t1 = 3 + 4} a \texttt{\$t1 = 7}.
  \item Propagation reemplaza \texttt{\$t2 = \$t1 * 2} por
        \texttt{\$t2 = 7 * 2}.
  \item Folding (siguiente iteracion) pliega a \texttt{\$t2 = 14}.
  \item Propagation reemplaza \texttt{c = \$t2} por \texttt{c = 14}.
  \item Dead-code elimina \texttt{\$t1} y \texttt{\$t2} muertos.
  \item Jump threading limpia etiquetas huerfanas si quedan.
\end{enumerate}
Por eso iteramos hasta punto fijo en lugar de una sola pasada.
\end{caja}

\section{Pruebas y QA: \texttt{qa\_rubrica.py}}

\begin{rubrica}{Cobertura contra la rubrica}
Suite de \textbf{42 casos automaticos} cubriendo los cinco criterios
tecnicos. Resultado actual: \textbf{42/42 PASS (100\%)}.
\end{rubrica}

\textbf{Desglose:}

\begin{lstlisting}[style=tac]
Criterio 1. Errores semanticos      20/20 casos PASS  (100%)
Criterio 2. Reporte HTML             2/2 casos PASS  (100%)
Criterio 3. TAC                     13/13 casos PASS  (100%)
Criterio 4. Optimizacion             5/5 casos PASS  (100%)
Criterio 5. Funcionamiento           2/2 casos PASS  (100%)
                                    -----
                                    42/42 PASS
\end{lstlisting}

\textbf{Ejemplos de los casos cubiertos (criterio 1):}
\begin{itemize}[leftmargin=*,nosep]
  \item Variable no declarada (uso simple y en asignacion).
  \item Variable duplicada en mismo ambito.
  \item Asignaciones int=cadena, int=decimal, bool=int.
  \item Widening decimal=entero permitido.
  \item Condicion de si/mientras/para no booleana.
  \item Division y modulo por cero literal.
  \item Suma entre booleano y entero.
  \item Comparacion entre tipos incompatibles.
  \item Operacion logica entre no booleanos.
  \item Operadores \texttt{-} y \texttt{!} con tipos invalidos.
\end{itemize}

\section{Cierre: cosas para senalar al presentar}

\begin{enumerate}[leftmargin=*,nosep]
  \item Mostrar el arbol de carpetas. Enfatizar que
        \texttt{intermedio.py} es la frontera.
  \item Hacer la \textbf{pausa teorica} sobre los tres tipos de IR.
        Escribir el ejemplo \texttt{(a+b)*2} en pizarra.
  \item Mostrar TAC original y TAC optimizado lado a lado. Enfatizar
        \emph{cuantos cuadruplos se redujeron} y \emph{que pasada lo
        logro}.
  \item Demostrar la traza del optimizador (iter, pasada, antes,
        despues, delta).
  \item Cargar un programa con errores semanticos en la UI y mostrar
        que \emph{el TAC no se genera} hasta arreglarlos.
  \item Hacer click en "REPORTES HTML" y abrir el reporte de errores
        semanticos en el navegador. Resaltar linea y columna.
  \item Mostrar la corrida de \texttt{python qa\_rubrica.py} con los
        42/42 PASS.
  \item Cerrar con la idea de \textbf{punto fijo}: las pasadas se
        alimentan entre si, por eso iteramos.
\end{enumerate}

\end{document}
